\section{Pastoral Use Without Turning a Reading into a Fact}

The daylight image is powerful precisely because it can be overstated. A
preacher may say that the chapter, read continuously, places Abram before stars
he cannot ordinarily see. A preacher should not say that Scripture explicitly
teaches that the stars were invisible, that Abraham was tested by daylight, or
that the tradition unanimously read the episode that way. Those are three
different claims, and the sources establish none of them.

The image also cannot carry counsel by itself. Trust in providence does not
require ignoring evidence, refusing prudent action, or treating uncertainty as
revelation. Abraham asks questions elsewhere in the same chapter; the covenant
promise does not abolish inquiry. The legitimate spiritual use is therefore
analogical and conditional: divine fidelity can exceed what is presently
visible. That proposition is Christianly intelligible without making a disputed
chronology into dogma.

For catechesis, the most useful lesson may be methodological. A vivid insight
should send the reader back to the whole passage, then outward to its versions
and reception. The result may strengthen the insight, narrow it, or retire it.
Here it survives, but only with its condition printed beside it.

